\chapter{Sức khỏe và thân thể}
\markboth{Sức khỏe và thân thể}{Health and Body}

\section{Người đàn ông chỉ đi khám khi cơn đau đã kéo dài nhiều tháng.}
\begin{truyen}{Người đàn ông chỉ đi khám khi cơn đau đã kéo dài nhiều tháng.}{Ignoring signs costs more later.}
\chuhoa{A}{nh sợ tốn thời gian, sợ nghe tin xấu, nên chọn cách chịu đựng. Đến khi bệnh nặng hơn, anh mới hiểu có những điều càng né càng đắt, cả bằng tiền lẫn bằng phần đời khỏe mạnh bị mất đi.}
\textit{(He feared wasting time and hearing bad news, so he chose to endure the pain. When the illness worsened, he learned that some things become more expensive the longer we avoid them, in both money and healthy years.)}
\end{truyen}
\begin{baihoc}
Trốn sự thật của cơ thể không làm cơ thể nương tay với ta.
\textit{(Avoiding the truth about your body does not make your body gentler with you.)}
\end{baihoc}
\begin{ghinhoanh}
\textbf{Short English to remember:}

Ignoring signs costs more later.

Listen before pain gets louder.
\end{ghinhoanh}
\begin{tuhocnhanh}
1. Có tín hiệu nào của cơ thể em đang coi nhẹ? \textit{(What signal from your body are you currently ignoring?)}

2. Em có thể làm gì để chăm sức khỏe chủ động hơn? \textit{(What can you do to care for your health more proactively?)}
\end{tuhocnhanh}
\ngancach

\section{Người trẻ đổi bữa ăn thành đồ uống cho nhanh.}
\begin{truyen}{Người trẻ đổi bữa ăn thành đồ uống cho nhanh.}{Convenience is not nourishment.}
\chuhoa{N}{hịp sống gấp khiến cô quen uống thứ tiện lợi thay cho bữa cơm tử tế. Vài năm sau, cơ thể mệt mỏi và tâm trạng dễ sụp xuống, cô mới hiểu không có kiểu ăn qua loa nào mà hậu quả cũng qua loa.}
\textit{(Her fast-paced life made her replace real meals with convenient drinks. Years later, when her body was tired and her mood unstable, she learned that careless eating never brings small consequences.)}
\end{truyen}
\begin{baihoc}
Thân thể cần được nuôi bằng sự chăm sóc đều đặn, không chỉ bằng thứ giúp ta sống sót qua ngày.
\textit{(The body must be nourished by steady care, not merely by whatever helps us survive the day.)}
\end{baihoc}
\begin{ghinhoanh}
\textbf{Short English to remember:}

Convenience is not nourishment.

Quick food can bring slow damage.
\end{ghinhoanh}
\begin{tuhocnhanh}
1. Em đang ăn để khỏe hay chỉ ăn để đỡ đói? \textit{(Are you eating to stay healthy or only to stop hunger?)}

2. Một thay đổi nhỏ nào trong bữa ăn em làm được ngay? \textit{(What small change can you make in your meals right away?)}
\end{tuhocnhanh}
\ngancach

\section{Một nhân viên xem kiệt sức là biểu hiện của chăm chỉ.}
\begin{truyen}{Một nhân viên xem kiệt sức là biểu hiện của chăm chỉ.}{Burnout is not a medal.}
\chuhoa{A}{nh tự hào vì lúc nào cũng mệt, coi đó như bằng chứng của nỗ lực. Cho đến khi đầu óc rỗng đi, tính khí xấu đi, và công việc giảm hẳn chất lượng, anh mới hiểu kiệt sức không phải huy chương, mà là hóa đơn.}
\textit{(He was proud of always being exhausted and treated it as proof of effort. When his mind emptied out, his temper worsened, and his work quality dropped, he realized burnout was not a medal but a bill.)}
\end{truyen}
\begin{baihoc}
Làm việc chăm không đồng nghĩa với làm việc đến mức tự phá mình.
\textit{(Working hard is not the same as working until you destroy yourself.)}
\end{baihoc}
\begin{ghinhoanh}
\textbf{Short English to remember:}

Burnout is not a medal.

Rest protects your work too.
\end{ghinhoanh}
\begin{tuhocnhanh}
1. Em có đang tự hào về những điều thật ra đang làm hại mình không? \textit{(Are you proud of things that are actually harming you?)}

2. Dấu hiệu mệt mỏi nào em đang xem nhẹ? \textit{(Which signs of exhaustion are you minimizing?)}
\end{tuhocnhanh}
\ngancach

\section{Cô gái chỉ thấy quý đôi chân sau một chấn thương.}
\begin{truyen}{Cô gái chỉ thấy quý đôi chân sau một chấn thương.}{We notice gifts after loss.}
\chuhoa{M}{ỗi ngày trước đó, cô đi lại bình thường mà không hề nghĩ ngợi. Sau vài tuần phải chống nạng, cô mới hiểu thứ ta sử dụng tự nhiên nhất lại thường là thứ ta biết ơn muộn nhất.}
\textit{(Before the injury, she walked normally without a thought. After weeks on crutches, she understood that what we use most naturally is often what we appreciate latest.)}
\end{truyen}
\begin{baihoc}
Biết ơn thân thể khi nó còn khỏe là một kiểu khôn ngoan hiếm gặp.
\textit{(Being grateful for your body while it is still healthy is a rare kind of wisdom.)}
\end{baihoc}
\begin{ghinhoanh}
\textbf{Short English to remember:}

We notice gifts after loss.

Thank your body before pain teaches you.
\end{ghinhoanh}
\begin{tuhocnhanh}
1. Bộ phận nào của cơ thể em đang được dùng mỗi ngày mà chưa từng được biết ơn? \textit{(Which part of your body do you use every day without gratitude?)}

2. Hôm nay em có thể chăm cơ thể mình bằng việc gì cụ thể? \textit{(How can you care for your body in a concrete way today?)}
\end{tuhocnhanh}
\ngancach

\section{Một người nghĩ tinh thần xuống dốc là do mình yếu đuối.}
\begin{truyen}{Một người nghĩ tinh thần xuống dốc là do mình yếu đuối.}{The mind also needs care.}
\chuhoa{A}{nh trách mình vì sao không vui nổi, vì sao dễ mệt và dễ cáu. Sau khi chịu lắng nghe chính mình và tìm hỗ trợ, anh mới hiểu sức khỏe tinh thần cũng là sức khỏe, không phải chuyện nên xấu hổ.}
\textit{(He blamed himself for not feeling happy, for being tired and irritable. After finally listening to himself and seeking support, he learned that mental health is also health, not something to be ashamed of.)}
\end{truyen}
\begin{baihoc}
Con người thường hiểu muộn rằng một tâm trí mệt mỏi cũng cần được chữa lành như một cơ thể đau.
\textit{(People often understand too late that a tired mind also needs healing, just like a hurting body.)}
\end{baihoc}
\begin{ghinhoanh}
\textbf{Short English to remember:}

The mind also needs care.

Asking for help is not weakness.
\end{ghinhoanh}
\begin{tuhocnhanh}
1. Dạo này tinh thần em đang thực sự thế nào? \textit{(How is your mental state really doing lately?)}

2. Em có thể tìm sự hỗ trợ từ ai nếu cần? \textit{(Who can you turn to for support if needed?)}
\end{tuhocnhanh}
\ngancach